\clearpage
\VerdierPhysicalPageAnchor{0108}%
\begingroup
\small
\setlength{\abovedisplayskip}{4pt plus 2pt minus 2pt}%
\setlength{\belowdisplayskip}{4pt plus 2pt minus 2pt}%

\VerdierSemanticAnchor{ch1-app-signs}
\section*{Appendice : Commentaires sur le choix des signes}
\addcontentsline{toc}{section}{Appendice : Commentaires sur le choix des signes}

\VerdierSemanticAnchor{ch1-app-signs-cone}
Le foncteur cône est déterminé à isomorphisme unique près par ses propriétés
d'adjonction (3.1.3). Pour pouvoir associer facilement à un morphisme de
$\operatorname{tr}_1(C)$ un triangle de $\operatorname{comp}(C)$, nous avons
choisi pour foncteur de translation le foncteur « cône du morphisme
$X\longrightarrow 0$ » ((3.1.7) et (3.2.2.7)), \emph{i.e.} décalage des
degrés et changement de signe de la différentielle.

\VerdierSemanticAnchor{ch1-app-signs-strict}
Soit
$F\colon\prod_{i\in[n]}C_i\longrightarrow C$ un foncteur multiadditif. La
catégorie $\prod_{i\in[n]}\operatorname{comp}(C_i)$ est munie naturellement
d'une structure de $\mathbb Z[n]$-catégorie. Si on veut que le foncteur
\[
  F^{[n]}\colon\prod_{i\in[n]}\operatorname{comp}(C_i)
  \longrightarrow\operatorname{comp}^{n}(C)
\]
soit un $\mathbb Z[n]$-foncteur strict, on est conduit au choix des
conventions que nous avons fait pour les translations : décalage des degrés
dans une direction et changement de signe de la différentielle
correspondante.

\VerdierSemanticAnchor{ch1-app-signs-simple-complex}
Nous avons pris les conventions de [1] pour le choix du foncteur complexe
simple associé (à ceci près que, dans nos complexes multiples, les
différentielles commutent entre elles). Ceci permet de définir le foncteur
$K_{\Pi}F\colon\prod_i K(C_i)\longrightarrow K(C)$. Il se trouve alors que
$K_{\Pi}F(X_1,\ldots,X_i[1],\ldots,X_n)$ est différent de
$K_{\Pi}F(X_1,\ldots,X_i,\ldots,X_n)[1]$. Pour pouvoir néanmoins affirmer que
le foncteur $K_{\Pi}F$ transforme les triangles distingués en triangles
distingués, on est amené à définir les isomorphismes :

\VerdierSemanticAnchor{ch1-app-signs-shift-isomorphism}
\[
  p_i\colon K_{\Pi}F(X_1,\ldots,X_i[1],\ldots,X_n)
  \xrightarrow{\ \sim\ }
  K_{\Pi}F(X_1,\ldots,X_i,\ldots,X_n)[1]
\]
(démonstration de (3.4.3)).

\VerdierSemanticAnchor{ch1-app-signs-anticommutativity}
Il se trouve alors que, pour $i\ne j$, on a des diagrammes
anticommutatifs :

\VerdierSemanticAnchor{ch1-app-signs-diagram}
\begin{center}
  \resizebox{\textwidth}{!}{%
    $\begin{tikzcd}[ampersand replacement=\&,column sep=3.2em,row sep=4.3em]
      K_{\Pi}F(X_1,\ldots,X_i[1],\ldots,X_j[1],\ldots)
        \arrow[r,"p_j"]
        \arrow[d,"p_i"']
        \& K_{\Pi}F(X_1,\ldots,X_i[1],\ldots,X_j,\ldots)[1]
          \arrow[d,"p_i"] \\
      K_{\Pi}F(X_1,\ldots,X_i,\ldots,X_j[1],\ldots)[1]
        \arrow[r,"p_j"]
        \& K_{\Pi}F(X_1,\ldots,X_i,\ldots,X_j,\ldots)[2].
    \end{tikzcd}$%
  }
\end{center}

\VerdierSemanticAnchor{ch1-app-signs-koszul-choice}
On a alors montré (1.5.2) comment prolonger de façon cohérente ces
isomorphismes en isomorphismes $p(\sigma)$ définis pour tout
$\sigma\in\mathbb Z[n]$. Il intervient ici un choix arbitraire. Il consiste à
choisir un cocycle dans une classe de cohomologie qui elle est bien
déterminée : la classe de Koszul (1.6.3). Pour nous conformer aux règles
communément admises (par exemple [1]), nous avons choisi de prendre dans tous
les cas le cocycle de Koszul (1.6.6).
\endgroup
